% =====================================================================================
%  Chapter 6 — Conclusion
%  Self-contained. Nothing to \input. Requires thesis/references.bib.
%
%  DEFINES sec:conc, the last dangling reference in the document (from background.tex).
%
%  NOVELTY LANGUAGE IS CALIBRATED AND SHOULD STAY THAT WAY. Three claims here are
%  deliberately narrower than they could be, and each is narrow for a reason:
%
%   1. "in the SafeLIBERO setting studied here" -- the transfer claim is empirical and
%      setting-bound, not general.
%   2. "not a transfer of the CBF's formal safety guarantee" -- Cosner et al. give conditions
%      under which a guarantee IS inherited; none holds here (Section 5.3 names which).
%   3. "the conjunction" -- learning safety from CBF supervision is NOT new, and distilling a
%      corrective signal into a VLA is NOT new (ROAD-VLA). Only the combination plus the
%      teacher-source boundary is.
%
%  Each of these would be easy to strengthen while editing and each would then be false.
% =====================================================================================

\chapter{Conclusion}\label{ch:conclusion}
\label{sec:conc}

This thesis asked whether the collision avoidance supplied by an inference-time control-barrier
function shield can be amortised into a vision-language-action policy, so that the shield need not
run at deployment. In the SafeLIBERO setting studied here the answer is yes, but only under a
specific form of supervision. Distilling $\pi_{0.5}$ on rollouts corrected by its own CBF shield
reduced collision episodes from 82.5\% to 19.2\% \textbf{without running the shield at inference},
while increasing task success from 58.3\% to 82.5\% on held-out initial states. The resulting
policy achieved a collision-avoidance rate of 80.8\%, approaching the shielded policy's 86.7\%, and
one distillation round was sufficient.

This extends the question posed by the work that introduced the benchmark~\cite{vlsa2026}. That
work improves $\pi_{0.5}$ on SafeLIBERO by \emph{retaining} a perception-grounded constraint layer
at inference; the present work asks what remains when its corrective behaviour is instead used as
training data and the layer is removed. What is demonstrated is an empirical reduction in
collisions, not a transfer of the CBF's formal safety guarantee: the student is no longer protected
by the online constraint satisfaction that establishes forward invariance~\cite{ames2017cbf}.
Conditions under which a learned controller does inherit input-to-state safety from a CBF-based
expert exist~\cite{cosner2022il_cbf}, but they require a robust expert, known control-affine
dynamics and explicit learning-theoretic assumptions, none of which this experiment establishes
(Section~\ref{sec:disc:limitations}).

\section*{The Boundary Condition}

The central boundary concerns the source and location of correction. A safe-by-construction,
geometry-privileged sampling-based planner did not improve the same student under a comparison
matched for gradient steps, hyperparameters and evaluation. Conversely, filtering only for
successful episodes also produced no measurable gain. Taken together, these controls support a
hypothesis about the \emph{states} supervision covers rather than about who produces it:
transferable supervision must correct errors on, or near, the states the learner itself reaches
when it begins to fail. Access to successful trajectories, even trajectories that are safe by
construction, is insufficient when they are drawn from a distribution the student does not visit.

The distinction matters because the stronger claim --- that a foreign controller cannot teach
safety --- is false, and the literature already shows it. SAFE-GIL trains behaviour cloning on
demonstrations from an MPC or PID expert, neither of which is the learner, and obtains
substantially safer policies~\cite{ciftci2025safegil}. What makes those demonstrations work is that
a reachability-computed disturbance is injected during collection, so the fixed expert is forced to
demonstrate recovery from the safety-critical states a learner's errors would produce; their
ablation confirms this targeting is the active ingredient, since undirected noise does not
reproduce the benefit. Their expert is privileged and foreign, and the demonstrations still
transfer. The planner tested here differs precisely in that its trajectories were generated from
its own initial conditions and were never conditioned on where $\pi_{0.5}$ goes wrong, and
Figure~\ref{fig:state_coverage} measures the consequence: at matched sampling it covers the
policy's own states substantially less well than the shielded rollouts do.

IntervenGen constructs the same idea independently, rolling out the \emph{current} policy so that
recorded mistakes are its genuine failures, then transforming a \emph{human} recovery segment onto
each state reached~\cite{hoque2024intervengen}. Its recovery actions, like SAFE-GIL's, come from
outside the policy, while the states they annotate come from the policy's own execution. Two
methods, with different teachers and different means of reaching the failure states, therefore
locate the active ingredient in the same place, and neither supports the stronger reading that the
teacher must be the policy itself.

The remaining alignment is with policy-proximal self-distillation~\cite{wang2026roadvla}, whose
design principle anticipates this result without testing it. That work constructs its teacher by
perturbing the student's own action-token logits, and derives a policy-improvement bound holding
only while the teacher stays proximal to the current policy. The relationship is a gap rather than
a disagreement: every teacher it rejects is a \emph{textual} one and every teacher it accepts is
derived from the student itself, so no foreign low-level controller appears anywhere in that work.
Its proximity principle predicts the planner result reported here, while the experiment that would
confirm it --- a privileged controller supplying well-formed low-level actions from outside the
policy --- is the one this thesis contributes.

Two features of the present result are not anticipated by that literature. Where SAFE-GIL reports a
safety--performance tradeoff, its own ablation showing task cost rising monotonically as more of
the training data moves into unsafe states, the distillation reported here improved collision
avoidance and task success \emph{together}. And where SAFE-GIL must synthesise the safety-critical
states adversarially, because its expert would not otherwise visit them, the shield supplies
corrections at states the policy reaches unprompted, so no disturbance model is required. Both
differences follow from correcting a policy on its own rollouts rather than steering a foreign
expert into the policy's likely failures.

\section*{What Bounds the Claim}

The remaining results sharpen, rather than universalise, that conclusion.

The distilled student avoided collisions involving arm links more effectively than the shield that
supervised it, indicating that trajectory imitation alone does not fully describe the learned
behaviour. That comparison is possible only because the barrier used here constrains links three
through seven and the hand, where the formulation it is measured against constrains the
end-effector alone by its authors' own statement~\cite{vlsa2026}. The effect is suggestive rather
than established --- the relevant Wilson intervals overlap (Section~\ref{sec:res:scope}).

Scalar-reward reinforcement learning, CBF-guided sampling, language-specified safety and
best-of-$N$ action selection did not provide an effective substitute under the implementations and
budgets tested. Each negative is scoped rather than general. The reinforcement-learning result
concerns the reward design tested and not safe reinforcement learning as a class: neither
constrained-MDP~\cite{safevla} nor model-based world-model approaches~\cite{safedojo} are
scalar-reward methods, and the latter is evaluated on this same benchmark. An independent study of
CBF-guided reinforcement learning reaches the same conclusion about the reward
channel~\cite{guerrier2025guardrails}, and supplies the control that makes the present result
non-trivial: training under an active shield and then removing it leaves their agent colliding, so
shield exposure alone does not transfer safety and the imitation objective is what does the work.
Best-of-$N$ likewise found little useful variation among samples drawn from the same state, which
supports a state-dominant account of collision risk \textbf{in this policy and benchmark} rather
than a general claim about generative policies.

\section*{What Is New, Stated Precisely}

Two neighbouring claims would be false and are not made. Learning safety from CBF-related
supervision is not new; it has been treated theoretically~\cite{cosner2022il_cbf}, and learning
from CBF demonstrations in simulation predates that work. Distilling a corrective signal into a
vision-language-action policy is not new either~\cite{wang2026roadvla}.

What is new is the conjunction: a control-theoretic safety filter, distilled from a VLA's own
shielded rollouts, evaluated with the filter \emph{removed}, in the setting where runtime shielding
is the current state of the art --- together with the boundary condition that a geometry-privileged
planner, safe by construction and supplying well-formed low-level actions, does not transfer under
a matched comparison. To this is added a measurement rather than a method: within a given state,
the safety outcomes of a VLA's sampled action candidates are near-perfectly correlated, which is
why action-level rejection sampling cannot help where episode-level selection can
(Section~\ref{sec:res:channels}).

The practical conclusion is therefore narrower than ``safety can be learned''. A runtime filter can
teach a policy the regularities that it observes and corrects, and policy-proximal corrections can
amortise a substantial part of that benefit without the per-step optimisation that comparable
learned-barrier methods retain at inference~\cite{meng2024conbat,sun2025lpb}. But the learned
policy is only as broad as the states, hazards and geometry represented in those corrections.
Runtime shielding remains necessary outside that validated envelope. Distillation reduces
dependence on the shield; it does not remove the need for a safety case.
